% Generated draft: requires mathematical and source review.
\documentclass[11pt]{article}
\usepackage[margin=25mm]{geometry}
\usepackage{fontspec}
\setmainfont{Latin Modern Roman}
\usepackage{amsmath,amssymb,mathtools,unicode-math}
\setmathfont{Latin Modern Math}
\usepackage{enumitem,needspace,xurl,hyperref}
\setlength{\parindent}{0pt}
\setlength{\parskip}{5pt}
\setlength{\emergencystretch}{3em}
\newcommand{\N}{\mathbb{N}}
\newcommand{\Z}{\mathbb{Z}}
\newcommand{\Q}{\mathbb{Q}}
\newcommand{\R}{\mathbb{R}}
\newcommand{\C}{\mathbb{C}}
\newcommand{\F}{\mathbb{F}}
\newcommand{\A}{\mathbb{A}}
\newcommand{\PP}{\mathbb{P}}
\newcommand{\E}{\mathbb{E}}
\newcommand{\eps}{\varepsilon}
\newcommand{\dd}{\,\mathrm{d}}
\begin{document}
\section*{UnsolvedMath ID: 6800014}
\textbf{Source problem code:} catalogue entry
\subsection{Definitions and mathematical statement}
Let \(G\) be a solvable Lie group with Lie algebra \(\mathfrak g\). Left-invariant Riemannian metrics are identified with positive definite inner products on \(\mathfrak g\), given their usual finite dimensional topology. Let \(\operatorname{Ric}_g\) be the Ricci endomorphism, defined by \(g(\operatorname{Ric}_g u,v)=\operatorname{Ricci}_g(u,v)\), and set
\[
 \operatorname{Scal}(g)=\operatorname{tr}(\operatorname{Ric}_g),\qquad
 |\operatorname{Ric}(g)|^2=\operatorname{tr}(\operatorname{Ric}_g^2),\qquad
 F(g)=\frac{\operatorname{Scal}(g)^2}{\operatorname{tr}(\operatorname{Ric}_g^2)}.
\]
The formula is defined only where \(\operatorname{Ric}_g\ne0\). On this domain, a local maximum is a metric \(g\) with a neighborhood \(U\) in which \(F(h)\leq F(g)\) for every admissible \(h\in U\). This is a nonstrict local maximum; no quotient by scaling or automorphisms is necessary for that definition.

A solvsoliton, in the usual algebraic sense, is a left-invariant metric satisfying
\[
 \operatorname{Ric}_g=cI+D\qquad(c\in\R,\ D\in\operatorname{Der}(\mathfrak g)),
\]
where \(D[u,v]=[Du,v]+[u,Dv]\) for all \(u,v\in\mathfrak g\). The question is whether every local maximum of \(F\) on the left-invariant metrics on a fixed solvable group is a solvsoliton. In particular, the requested implication is
\[
 g\text{ locally maximizes }F\quad\Longrightarrow\quad
 \operatorname{Ric}_g=cI+D.
\]
It is not a request to classify only global maximizers. The source does not specify how to treat Ricci-flat metrics, including the identically undefined functional on an abelian group. Restricting to its nonzero-denominator domain is necessary for the displayed formula and is marked for review.
\subsection{Short English statement}
Must every local maximizer of scalar curvature squared divided by squared Ricci norm, among left-invariant metrics on a solvable Lie group, be an algebraic Ricci soliton? The functional's undefined Ricci-flat locus needs a convention.
\subsection{Sources}
\begin{itemize}
\item Local UnsolvedMath registry, permanent ID 6800014, source code catalogue entry; attributed to Jorge Lauret, C\'ordoba, Argentina; proposed year 2018; update 2026-08-06.
\item Morgan and Pansu, \emph{A List of Open Problems in Differential Geometry} (2018), Question environment 14, ``Ricci pinching on solvable Lie groups'': \url{https://www.imo.universite-paris-saclay.fr/~pansu/problems_MTDG.tex}. Registry access date 2026-07-29; source-TeX extraction; public source on a coauthor's publication page; current status and redistribution terms require review.
\item The supplied summary reports global-maximization results in unimodular, almost abelian, codimension one nilradical, and certain abelian nilradical cases, while describing the general local-maximum statement as unresolved. No exact further reference or hypotheses are supplied.
\end{itemize}
\end{document}
